\chapter[Estudos de simulação]{Estudos de simulação}\label{sec-estudos-simulacao}
A fim de avaliar o comportamento empírico dos estimadores propostos, conduzimos dois estudos de simulação principais. O primeiro compara a nossa proposta baseada em ondaletas --- Seção \ref{sec-estimacao-fate} em suas versões \emph{offline} e \emph{online} --- com o estimador de \emph{splines} utilizado em \citeonline{ecker_causal_2024} (Seção \ref{sec-inferencia-causal-funcional}). Como já discutido, ao adotarmos o mesmo processo gerador de dados (PGD) e o cenário de intervenção única dos autores, tornamos a exigência de causalidade operacional (respeito à seta do tempo) optativa. Esse desenho nos permite comparar diretamente a eficiência da suavização e a identificação do efeito entre os métodos, estabelecendo o estimador \emph{offline} como um teto teórico de desempenho para a versão \emph{online}.\footnote{Uma vez que o cenário não requer o respeito à causalidade operacional, o estimador \emph{offline} deve desempenhar melhor que o \emph{online} em virtude de ter acesso a uma quantidade maior de informação.} A segunda simulação, por sua vez, foca exclusivamente em nosso estimador \emph{online} (Seção \ref{sec-estimador-online}; tendo em vista que não encontramos abordagens alternativas para estimar o \(\mathrm{FATE}\) sequencialmente na literatura), mensurando a latência computacional do processamento ponto a ponto e a precisão da estimativa quando as observações de cada indivíduo são disponibilizadas sequencialmente, sem que o estimador tenha acesso a informações futuras --- por óbvio, o desempenho estatístico do estimador \emph{online} não deve mudar significativamente de um desenho experimental para outro. A única diferença entre os estimadores é que, na primeira simulação, utilizamos a função \texttt{denoise\_signal\_causal}, enquanto na segunda, emulamos os dados chegando sequencialmente com \texttt{new\_wavelet\_stream}. Ambas as funções estão disponíveis no ecossistema do \texttt{rLifting}.

Em nosso desenho, buscamos mimetizar o PGD de \citeonline{ecker_causal_2024} para garantir uma base de comparação mais sólida. Geramos \(P = 9\) covariáveis pré-tratamento amostradas de distribuições Normal, Poisson e Binomial, cujos parâmetros foram calibrados com base nas estatísticas descritivas reportadas no trabalho original (Tabela \ref{tab:parametros_pgd}). A atribuição de tratamento segue um modelo logístico com viés de seleção induzido, e as propensões foram truncadas no intervalo \([0.001, 0.999]\) para assegurar a positividade condicional. O desfecho funcional é gerado a partir de um modelo concorrente no qual o efeito de tratamento é heterogêneo, dependendo das covariáveis observadas. Distribuímos o efeito total uniformemente entre o intercepto e as \(9\) covariáveis, de modo que cada componente contribua igualmente para o \(\mathrm{FATE}\) verdadeiro através da média populacional \(\bar{\mathbf{X}}\).

A fim de exemplificar os ganhos de nossos estimadores, quando comparado ao PGD de \citeonline{ecker_causal_2024}, ampliamos o conjunto de formas funcionais avaliadas. Consideramos funções de efeito com os formatos: \emph{sigmoide}, \emph{bumps}, \emph{doppler} e \emph{escada}. Esperávamos, a priori, que o estimador original, baseado em \emph{splines}, desempenhasse mal nos últimos três cenários supracitados.

Em ambas as simulações cruzamos tamanhos de amostra \(N \in \{200, 500, 2000\}\) com resoluções temporais \(T \in \{128, 256, 512\}\). Para o pré-processamento via ondaletas, a família mãe foi selecionada conforme o cenário (utilizamos filtros curtos para o caso \emph{online}, conforme nossa discussão metodológica, e para o caso \emph{offline}, por comparabilidade; e.g., Haar para o cenário Escada e Daubechies-2 para Bumps), enquanto para os \emph{splines} utilizamos o padrão de \(50\) funções de base em todos os casos (otimizamos o procedimento via validação cruzada generalizada --- o que é de praxe e, inclusive, é feito por \citeonline{ecker_causal_2024}, apesar de teoricamente afetar a normalidade esperada do estimador \(\hat{\theta}(t)\) resultante). Onde aplicável, a janela de processamento causal foi fixada em \(W = 64\) para resoluções menores e \(W = 128\) para \(T = 512\), permitindo um equilíbrio entre a estabilidade da reconstrução e o respeito à seta do tempo. 

Obtemos as estimativas intervalares para os métodos baseados em ondaletas através do procedimento descrito na Seção \ref{sec-bootstrap-ic}. Para as \emph{splines}, utilizamos as derivações analíticas de \citeonline{ecker_causal_2024}. Vale salientar, no entanto, que introduzimos uma escolha deliberadamente conservadora na calibração do limiar de cobertura \(u_\alpha(t)\). Em vez de utilizar a estrutura variável original, baseada no pacote \texttt{ffscb} \cite{lieblFastFairSimultaneous2023}, fixamos um limiar escalar constante correspondente ao valor mais conservador do referido pacote. O objetivo desta modificação era assegurar que eventuais falhas de cobertura do método original em cenários irregulares não decorressem de bandas excessivamente estreitas localmente. Estratégia que nos permite uma avaliação mais isolada do impacto da falta de regularidade na estimativa pontual. Essa escolha, é claro, inviabiliza a comparação das larguras de banda resultantes de nosso \emph{bootstrap}, para com o método de \citeonline{ecker_causal_2024} (no fim das contas, mesmo com nossas precauções, a comparação seria inócua, tendo em vista que a metodologia original não se mostrou estatisticamente confiável para cenários irregulares mesmo com larguras de banda artificialmente maiores).

\begin{table}[ht]
	\centering
	\caption{Especificação paramétrica do processo gerador de dados das simulações.}
	\label{tab:parametros_pgd}
	\begin{tabular}{llrr}
		\toprule
		Covariável & Distribuição & Escore de propensão & Efeito concorrente \\
		\midrule
		\(X_{i1}\) & \(\mathrm{N}(-2.3;\ 1.2^2)\) & \(-0.052\) & \(-0.02\) \\
		\(X_{i2}\) & \(\mathrm{N}(8.3;\ 4.1^2)\) & \(-0.001\) & \(0.5\exp(-2t)\) \\
		\(X_{i3}\) & \(\mathrm{Bin}(1;\ 0.019)\) & \(0.020\) & \(-0.1 + 0.2(t-0.5)^2\) \\
		\(X_{i4}\) & \(\mathrm{Bin}(1;\ 0.710)\) & \(-0.097\) & \(0.2t\) \\
		\(X_{i5}\) & \(\mathrm{Pois}(2.3)\) & \(-0.056\) & \(0.01\sin(10t)\) \\
		\(X_{i6}\) & \(\mathrm{N}(11.8;\ 0.9^2)\) & \(0.510\) & \(0.05\) \\
		\(X_{i7}\) & \(\mathrm{Bin}(1;\ 0.023)\) & \(0.568\) & \(-0.1\) \\
		\(X_{i8}\) & \(\mathrm{Bin}(1;\ 0.297)\) & \(0.674\) & \(0.1\exp(-t)\) \\
		\(X_{i9}\) & \(\mathrm{Bin}(1;\ 0.112)\) & \(0.851\) & \(0.15\exp(-t)\) \\
		\bottomrule
	\end{tabular}
\end{table}

Para avaliar o desempenho dos estimadores, consideramos métricas de precisão pontual e global, além de propriedades de cobertura das bandas de confiança simultâneas. Denotamos por \(\theta(t_k)\) o efeito populacional verdadeiro no instante \(t_k \equiv k\) (com \(k = 0, \dots, T-1\)) e por \(\hat{\theta}^{(r)}(t_k)\) a sua respectiva estimativa obtida na réplica \(r = 1, \dots, R\) da simulação (com \(R = 1000\)).

Iniciamos pela avaliação da precisão em cada instante de tempo, calculando o erro quadrático médio e o viés pontuais através da média sobre as \(R\) réplicas:
\begin{equation}
	\mathrm{RMSE}(t_k) = \sqrt{\frac{1}{R}\sum_{r=1}^{R} \bigl(\hat{\theta}^{(r)}(t_k) - \theta(t_k)\bigr)^2}, \quad \text{Viés}(t_k) = \frac{1}{R}\sum_{r=1}^{R} \bigl(\hat{\theta}^{(r)}(t_k) - \theta(t_k)\bigr)
\end{equation}
Essas medidas permitem visualizar a evolução do erro e a estabilidade da estimativa ao longo das trajetórias. Em seguida, definimos as métricas globais para cada réplica \(r\), agregadas sobre os \(T\) instantes de tempo:
\begin{description}
	\item[Raíz quadrada do erro quadrático médio (RMSE)] quantifica a precisão global da trajetória estimada em cada réplica:
	\begin{equation*}
		\\mathrm{RMSE}^{(r)} = \sqrt{\frac{1}{T}\sum_{k=0}^{T-1} \bigl(\hat{\theta}^{(r)}(t_k) - \theta(t_k)\bigr)^2}.
	\end{equation*}
	\item[Raíz quadrada do erro quadrático médio normalizado (NRMSE)] escalona o \text{RMSE} pela amplitude da trajetória verdadeira, facilitando a comparação entre cenários com diferentes magnitudes de efeito:
	\begin{equation*}
		\mathrm{NRMSE}^{(r)} = \frac{\text{RMSE}^{(r)}}{\max_{t_k} \theta(t_k) - \min_{t_k} \theta(t_k)}.
	\end{equation*}
	\item[Viés global] avalia o erro sistemático médio ao longo de toda a curva:
	\begin{equation*}
		\text{Viés}^{(r)} = \frac{1}{T}\sum_{k=0}^{T-1} \bigl(\hat{\theta}^{(r)}(t_k) - \theta(t_k)\bigr).
	\end{equation*}
	\item[Latência computacional] representa o tempo médio de processamento por ponto na grade temporal (exclusivo para a Simulação 2):
	\begin{equation*}
		\text{Latência}^{(r)} = \frac{1}{T}\sum_{k=0}^{T-1} \text{Tempo}(\text{Processamento ponto } k).
	\end{equation*}
	\item[Cobertura simultânea] é uma indicadora binária que assume \(1\) se, e somente se, a banda de confiança contém a trajetória verdadeira em todos os instantes \(t_k\), simultaneamente:
	\begin{equation*}
		\mathrm{CobSim}^{(r)} = \mathbb{I}\left( \forall k : \theta(t_k) \in [\hat{\theta}_{\text{low}}^{(r)}(t_k), \hat{\theta}_{\text{up}}^{(r)}(t_k)] \right).
	\end{equation*}
	\item[Largura média da banda de confiança], monitorada a fim de avaliar o custo informacional necessário para atingir o nível de confiança nominal.
	\begin{equation*}
		\(\mathrm{Largura}^{(r)} = \frac{1}{T}\sum_{k=0}^{T-1} \bigl(\hat{\theta}_{\text{up}}^{(r)}(t_k) - \hat{\theta}_{\text{low}}^{(r)}(t_k)\bigr)\)
	\end{equation*}

\end{description}

Deve-se observar, contudo, que a métrica de latência reportada representa a média de processamento por ponto na grade, obtida em ambiente controlado com hardware comercial padrão. Embora os valores absolutos possam variar conforme a infraestrutura, a estabilidade e a ordem de grandeza observada confirmam a viabilidade operacional do método para o regime sequencial pretendido.

A seguir, apresentamos os resultados obtidos. Para fins de clareza e síntese no corpo principal deste trabalho, optamos por exibir uma seleção representativa dos cenários (Sigmoid, Doppler e Escada), contrastando os extremos de resolução temporal (\(T = 128\) e \(T = 512\)) sob diferentes tamanhos de amostra. Esta escolha permite observar os padrões de convergência e as vantagens competitivas da metodologia proposta sem sobrecarregar a leitura com a totalidade das combinações simuladas. Os resultados exaustivos para todos os cenários e resoluções encontram-se detalhados no Apêndice \ref{ap-simulacoes-completas}.

Note que, para fins de clareza visual, as distribuições de erro (\(\mathrm{NRMSE}\)) e largura de banda nas figuras subsequentes são representadas pela mediana e pelos quartis de 25\% e 75\%, de modo a omitir valores atípicos que dificultariam a percepção visual do impacto das resoluções \(T\) e \(N\).

\section{Resultados}

Abaixo, apresentamos as tabelas gerenciais que condensam o desempenho dos métodos avaliados. No que tange à Simulação 1, focamos na comparação direta entre Ondaletas e Splines de Ecker para os cenários mais desafiadores.

\begin{table}[ht]
	\centering
	\caption{Desempenho estatístico comparado (Simulação 1) --- Resultados Selecionados.}
	\label{tab:resultados_sim1}
	\scriptsize
	\begin{tabular}{llccrrrr}
		\toprule
		Cenário & Método & \(N\) & \(T\) & \(\text{RMSE}\) & \(\text{NRMSE}\) & \(\text{Vício}\) & \(\text{CobSim}\) \\
\midrule
\multirow{18}{*}{Sigmoid} & Ond. (On) & 200 & 128 & 1,278 & 0,574 & 0,055 & 85,9\% \\
 & Ond. (Off) & 200 & 128 & 1,304 & 0,585 & 0,058 & 86,6\% \\
 & Ecker & 200 & 128 & 1,300 & 0,584 & 0,058 & 100,0\% \\
\cmidrule{2-8}
 & Ond. (On) & 200 & 512 & 1,199 & 0,538 & 0,106 & 88,9\% \\
 & Ond. (Off) & 200 & 512 & 1,206 & 0,541 & 0,107 & 89,5\% \\
 & Ecker & 200 & 512 & 1,205 & 0,541 & 0,107 & 100,0\% \\
\cmidrule{2-8}
 & Ond. (On) & 500 & 128 & 0,091 & 0,041 & 0,009 & 96,8\% \\
 & Ond. (Off) & 500 & 128 & 0,099 & 0,044 & 0,011 & 97,4\% \\
 & Ecker & 500 & 128 & 0,091 & 0,041 & 0,011 & 100,0\% \\
\cmidrule{2-8}
 & Ond. (On) & 500 & 512 & 0,102 & 0,046 & 0,007 & 96,1\% \\
 & Ond. (Off) & 500 & 512 & 0,103 & 0,046 & 0,008 & 96,9\% \\
 & Ecker & 500 & 512 & 0,102 & 0,046 & 0,008 & 100,0\% \\
\cmidrule{2-8}
 & Ond. (On) & 2000 & 128 & 0,022 & 0,010 & -0,001 & 96,1\% \\
 & Ond. (Off) & 2000 & 128 & 0,034 & 0,015 & 0,001 & 94,7\% \\
 & Ecker & 2000 & 128 & 0,022 & 0,010 & 0,001 & 100,0\% \\
\cmidrule{2-8}
 & Ond. (On) & 2000 & 512 & 0,021 & 0,010 & -0,001 & 94,6\% \\
 & Ond. (Off) & 2000 & 512 & 0,023 & 0,010 & 0,000 & 95,0\% \\
 & Ecker & 2000 & 512 & 0,021 & 0,010 & 0,000 & 100,0\% \\
\cmidrule{2-8}
\midrule
\multirow{18}{*}{Doppler} & Ond. (On) & 200 & 128 & 1,785 & 0,161 & 0,027 & 86,8\% \\
 & Ond. (Off) & 200 & 128 & 1,716 & 0,155 & 0,044 & 86,8\% \\
 & Ecker & 200 & 128 & 2,806 & 0,253 & 0,051 & 40,1\% \\
\cmidrule{2-8}
 & Ond. (On) & 200 & 512 & 1,498 & 0,135 & 0,071 & 87,1\% \\
 & Ond. (Off) & 200 & 512 & 1,253 & 0,113 & 0,078 & 87,1\% \\
 & Ecker & 200 & 512 & 2,337 & 0,210 & 0,079 & 40,3\% \\
\cmidrule{2-8}
 & Ond. (On) & 500 & 128 & 0,805 & 0,073 & -0,003 & 96,1\% \\
 & Ond. (Off) & 500 & 128 & 0,722 & 0,065 & 0,014 & 95,8\% \\
 & Ecker & 500 & 128 & 2,011 & 0,182 & 0,021 & 3,8\% \\
\cmidrule{2-8}
 & Ond. (On) & 500 & 512 & 0,467 & 0,042 & -0,001 & 96,2\% \\
 & Ond. (Off) & 500 & 512 & 0,136 & 0,012 & 0,006 & 96,8\% \\
 & Ecker & 500 & 512 & 1,493 & 0,134 & 0,007 & 3,0\% \\
\cmidrule{2-8}
 & Ond. (On) & 2000 & 128 & 0,736 & 0,066 & -0,025 & 94,9\% \\
 & Ond. (Off) & 2000 & 128 & 0,650 & 0,059 & -0,008 & 94,8\% \\
 & Ecker & 2000 & 128 & 1,960 & 0,177 & -0,001 & 0,0\% \\
\cmidrule{2-8}
 & Ond. (On) & 2000 & 512 & 0,406 & 0,037 & -0,008 & 96,2\% \\
 & Ond. (Off) & 2000 & 512 & 0,060 & 0,005 & -0,001 & 95,2\% \\
 & Ecker & 2000 & 512 & 1,445 & 0,130 & 0,000 & 0,0\% \\
\cmidrule{2-8}
\midrule
\multirow{18}{*}{Escada} & Ond. (On) & 200 & 128 & 1,235 & 0,046 & 0,084 & 87,2\% \\
 & Ond. (Off) & 200 & 128 & 1,257 & 0,047 & 0,077 & 88,0\% \\
 & Ecker & 200 & 128 & 2,283 & 0,086 & 0,077 & 40,8\% \\
\cmidrule{2-8}
 & Ond. (On) & 200 & 512 & 1,347 & 0,050 & 0,088 & 85,7\% \\
 & Ond. (Off) & 200 & 512 & 1,354 & 0,051 & 0,087 & 86,0\% \\
 & Ecker & 200 & 512 & 2,100 & 0,079 & 0,087 & 42,7\% \\
\cmidrule{2-8}
 & Ond. (On) & 500 & 128 & 0,125 & 0,005 & 0,005 & 97,7\% \\
 & Ond. (Off) & 500 & 128 & 0,127 & 0,005 & -0,003 & 97,2\% \\
 & Ecker & 500 & 128 & 1,427 & 0,053 & -0,003 & 3,4\% \\
\cmidrule{2-8}
 & Ond. (On) & 500 & 512 & 0,086 & 0,003 & 0,027 & 97,4\% \\
 & Ond. (Off) & 500 & 512 & 0,086 & 0,003 & 0,026 & 97,3\% \\
 & Ecker & 500 & 512 & 1,073 & 0,040 & 0,026 & 2,6\% \\
\cmidrule{2-8}
 & Ond. (On) & 2000 & 128 & 0,049 & 0,002 & 0,008 & 95,4\% \\
 & Ond. (Off) & 2000 & 128 & 0,050 & 0,002 & 0,000 & 95,9\% \\
 & Ecker & 2000 & 128 & 1,379 & 0,052 & 0,000 & 0,0\% \\
\cmidrule{2-8}
 & Ond. (On) & 2000 & 512 & 0,022 & 0,001 & 0,001 & 95,5\% \\
 & Ond. (Off) & 2000 & 512 & 0,022 & 0,001 & 0,000 & 95,7\% \\
 & Ecker & 2000 & 512 & 1,040 & 0,039 & 0,000 & 0,0\% \\
\cmidrule{2-8}
		\bottomrule
	\end{tabular}
\end{table}

\begin{table}[ht]
	\centering
	\caption{Eficiência computacional e confiabilidade (Simulação 2) --- Ondaletas Online.}
	\label{tab:resultados_sim2}
	\begin{tabular}{lcrrrrr}
		\toprule
		Cenário & \(T\) & \(N\) & \(\text{RMSE}\) & \(\text{NRMSE}\) & \(\text{Vício}\) & \(\text{Latência (µs)}\) \\
\midrule
\multirow{9}{*}{Bumps} & 128 & 200 & 1,391 & 0,022 & 0,244 & 3,7 \\
 & 128 & 500 & 0,180 & 0,003 & 0,009 & 3,9 \\
 & 128 & 2000 & 0,116 & 0,002 & -0,015 & 4,1 \\
\cmidrule{2-7}
 & 256 & 200 & 1,323 & 0,020 & 0,041 & 4,2 \\
 & 256 & 500 & 0,208 & 0,003 & 0,028 & 4,2 \\
 & 256 & 2000 & 0,122 & 0,002 & -0,021 & 5,1 \\
\cmidrule{2-7}
 & 512 & 200 & 1,432 & 0,021 & -0,058 & 5,3 \\
 & 512 & 500 & 0,101 & 0,002 & 0,017 & 5,4 \\
 & 512 & 2000 & 0,040 & 0,001 & -0,002 & 6,9 \\
\cmidrule{2-7}
\midrule
\multirow{9}{*}{Doppler} & 128 & 200 & 1,827 & 0,165 & 0,015 & 3,7 \\
 & 128 & 500 & 0,771 & 0,070 & -0,024 & 3,8 \\
 & 128 & 2000 & 0,736 & 0,066 & -0,024 & 4,1 \\
\cmidrule{2-7}
 & 256 & 200 & 1,937 & 0,175 & 0,048 & 4,2 \\
 & 256 & 500 & 0,952 & 0,086 & -0,049 & 4,2 \\
 & 256 & 2000 & 0,852 & 0,077 & -0,067 & 7,6 \\
\cmidrule{2-7}
 & 512 & 200 & 1,586 & 0,143 & 0,101 & 5,3 \\
 & 512 & 500 & 0,480 & 0,043 & 0,001 & 5,5 \\
 & 512 & 2000 & 0,405 & 0,036 & -0,008 & 8,7 \\
\cmidrule{2-7}
\midrule
\multirow{9}{*}{Heavisine} & 128 & 200 & 1,245 & 0,047 & 0,041 & 3,5 \\
 & 128 & 500 & 0,115 & 0,004 & 0,003 & 3,6 \\
 & 128 & 2000 & 0,049 & 0,002 & 0,008 & 3,9 \\
\cmidrule{2-7}
 & 256 & 200 & 1,457 & 0,055 & 0,119 & 4,0 \\
 & 256 & 500 & 0,114 & 0,004 & 0,001 & 4,0 \\
 & 256 & 2000 & 0,031 & 0,001 & 0,004 & 4,3 \\
\cmidrule{2-7}
 & 512 & 200 & 1,106 & 0,041 & 0,061 & 4,7 \\
 & 512 & 500 & 0,081 & 0,003 & 0,026 & 5,0 \\
 & 512 & 2000 & 0,023 & 0,001 & 0,001 & 5,2 \\
\cmidrule{2-7}
\midrule
\multirow{9}{*}{Sigmoid} & 128 & 200 & 1,315 & 0,591 & 0,046 & 3,6 \\
 & 128 & 500 & 0,093 & 0,042 & 0,006 & 3,6 \\
 & 128 & 2000 & 0,022 & 0,010 & -0,001 & 3,9 \\
\cmidrule{2-7}
 & 256 & 200 & 1,179 & 0,529 & 0,135 & 4,0 \\
 & 256 & 500 & 0,085 & 0,038 & 0,008 & 4,0 \\
 & 256 & 2000 & 0,022 & 0,010 & -0,002 & 4,7 \\
\cmidrule{2-7}
 & 512 & 200 & 1,295 & 0,582 & 0,083 & 4,7 \\
 & 512 & 500 & 0,105 & 0,047 & 0,028 & 5,0 \\
 & 512 & 2000 & 0,021 & 0,009 & 0,000 & 5,2 \\
\cmidrule{2-7}
		\bottomrule
	\end{tabular}
\end{table}
